% ભાગ: ગણો-વિધેયો-સંબંધો
% પ્રકરણ: અંકગણિતીકરણ
% વિભાગ: વિગતોની ચકાસણી
%
\documentclass[../../../include/open-logic-section]{subfiles}


\begin{document}

\olfileid{sfr}{arith}{check}
\olsection{ક્રમિત મંડળો અને ક્ષેત્રો}

આ આખા પ્રકરણમાં આપણે દાવો કર્યો કે અમુક વ્યાખ્યાઓનો વ્યવહાર
``જેવો હોવો જોઈએ તેવો'' છે. આ પ્રાવિધિક પરિશિષ્ટમાં તેનો અર્થ
સ્પષ્ટ કરીશું અને વ્યાખ્યાઓનો વ્યવહાર ``યોગ્ય રીતે'' થાય છે તે
(કેવી રીતે બતાવવું તેની રૂપરેખા) આપીશું.

\olref[int]{sec}માં આપણે $\Int$ પર સરવાળો અને ગુણાકાર વ્યાખ્યાયિત
કર્યા. વ્યાખ્યા મુજબ આ ક્રિયાઓ $\Int$ને આપણે ``ઇચ્છીએ તેવું''
સંરચના આપે છે એવું બતાવવું છે. ખાસ કરીને, અહીંની સંરચના સમક્રમી
મંડળની છે.

\begin{defn}
	\emph{સમક્રમી મંડળ} એવો ગણ $S$ છે જેમાં નિશ્ચિત ઘટકો $0$ અને $1$
	તથા ક્રિયાઓ $+$ અને $\times$ હોય અને જે નીચેના આઠ સૂત્રો સંતોષે:
	\begin{align*}
		\emph{સંગઠિતતા}&&a + (b+ c) & = (a + b) + c \\
		&& (a \times b) \times c & = a \times (b\times c)\\
		\emph{સમક્રમિતા}&&a + b &= b+ a  \\
		&&  a \times b&= b\times a\\
		\emph{એકમો}&&a + 0 &= a \\
		&& a \times 1 &= a\\
		\emph{યોજક વ્યસ્ત}&&(\exists b\in S)0&=a + b\\
		\emph{વિતરણાત્મકતા}&&a \times (b+ c ) &= (a \times b) + (a \times c)
	\end{align*}
	અહીં ગર્ભિત રીતે આ બધાં સૂત્રોના સર્વપરિમાણકો $S$ પૂરતા
	મર્યાદિત છે. એ પણ નોંધો કે અહીંના $0$ અને~$1$ નામના ઘટકો એ જ
	નામ ધરાવતી પ્રાકૃતિક સંખ્યાઓ હોવા જરૂરી નથી.
\end{defn}

તેથી પૂર્ણાંકો સમક્રમી મંડળ બનાવે છે કે નહીં તે ચકાસવા માટે આ આઠે
શરતો સંતોષાય છે એટલું જ ચકાસવાનું રહે છે. કોઈ શરત સ્થાપવી {મુશ્કેલ}
નથી, પણ એ કામ થોડું શ્રમસાધ્ય છે. દાખલા તરીકે, સરવાળાના કિસ્સામાં
\emph{સંગઠિતતા} આ રીતે સાબિત કરી શકાય:

\begin{proof}
$i, j, k \in \Int$ નિશ્ચિત કરો. ત્યારે એવા $a_1, b_1, a_2, b_2, a_3, b_3 \in
\Nat$ છે કે $i = \equivrep{a_1, b_1}{}$, $j =
\equivrep{a_2,b_2}{}$ અને $k = \equivrep{a_3, b_3}{}$. (વાંચવામાં
સરળતા રહે તે માટે આપણે ``$\equivrep{x, y}{\Intequiv}$''ને બદલે
``$\equivrep{x, y}{}$'' લખીએ છીએ; આ આખા વિભાગમાં એમ જ કરીશું.) હવે:
\begin{align*}
	i + (j + k) &= \equivrep{a_1, b_1}{}+(\equivrep{a_2, b_2}{} + \equivrep{a_3, b_3}{}) \\
	&= \equivrep{a_1,  b_1}{} + \equivrep{a_2+a_3, b_2+b_3}{}\\
	&= \equivrep{a_1 + (a_2 + a_3), b_1 + (b_2 + b_3)}{}\\
	&= \equivrep{(a_1 + a_2) + a_3, (b_1 + b_2) + b_3}{}\\
	&= \equivrep{a_1 + a_2, b_1 + b_2}{} + \equivrep{a_3, b_3}{}\\
	&= (\equivrep{a_1, b_1}{} + \equivrep{a_2, b_2}{}) + \equivrep{a_3, b_3}{}\\
	&= (i+j) + k
\end{align*}
અહીં $\Nat$ પરના સરવાળાના વ્યવહારનો આપણે નિઃસંકોચ ઉપયોગ કર્યો છે.
\end{proof}

એ જ રીતે \emph{યોજક વ્યસ્ત} આ રીતે સાબિત કરી શકાય:

\begin{proof}
$i \in \Int$ નિશ્ચિત કરો; ત્યારે કોઈ $a,b \in \Nat$ માટે $i =
\equivrep{a,b}{}$. $j = \equivrep{b,a}{} \in \Int$ લો. પ્રાકૃતિક
સંખ્યાઓના વ્યવહારનો ઉપયોગ કરતાં $(a+b) + 0 = 0 + (a+b)$, તેથી
વ્યાખ્યા મુજબ $\tuple{a+b, b+a} \sim_\Int \tuple{0,0}$, અને તેથી
$\equivrep{a+b, b+a}{} = \equivrep{0, 0}{} = 0_\Int$. હવે $i + j =
\equivrep{a,b}{}+\equivrep{b,a}{}=\equivrep{a+b, b+a}{}= \equivrep{0,
0}{} = 0_\Int$.
\end{proof}

અને \emph{વિતરણાત્મકતા}ની સાબિતી આ રહી:

\begin{proof}
ઉપરની જેમ $i = \equivrep{a_1, b_1}{}$, $j =
\equivrep{a_2,b_2}{}$ અને $k = \equivrep{a_3, b_3}{}$ નિશ્ચિત કરો. હવે:
\begin{align*}
	i \times (j + k)
	&= \equivrep{a_1, b_1}{} \times (\equivrep{a_2,b_2}{} + \equivrep{a_3, b_3}{})\\
	&= \equivrep{a_1, b_1}{} \times \equivrep{a_2 + a_3,b_2+b_3}{}\\
	&= \equivrep{a_1  (a_2 + a_3) + b_1  (b_2+b_3), a_1  (b_2 + b_3) + b_1 (a_2 + a_3)}{}\\
	&= \equivrep{a_1 a_2 + a_1a_3 + b_1 b_2+b_1b_3, a_1 b_2 + a_1b_3 + a_2b_1 + a_3b_1}{}\\
%		&= \equivrep{a_1 a_2 + b_1b_2 + a_1a_3 + b_1 b_2+b_1b_3, a_1 b_2 + a_1b_3 + a_2b_1 + a_3b_1}{}\\
	&= \equivrep{a_1a_2 + b_1b_2, a_1b_2 + a_2b_1}{} + \equivrep{a_1a_3 + b_1b_3, a_1b_3 + a_3b_1}{}\\
	&= (\equivrep{a_1, b_1}{} \times \equivrep{a_2,b_2}{}) + (\equivrep{a_1, b_1}{} \times  \equivrep{a_3, b_3}{})\\
	&= (i \times j) + (i \times k)
\end{align*}
\end{proof}

બાકીની પાંચ શરતોની સાબિતી કસરત તરીકે છોડી દઈએ છીએ. એ કરી લઈએ
પછી આપણે બતાવ્યું હશે કે $\Int$ સમક્રમી મંડળ બનાવે છે, એટલે કે સરવાળો
અને ગુણાકારનો (વ્યાખ્યાયિત કરેલો) વ્યવહાર જેવો હોવો જોઈએ તેવો છે.

\begin{prob}
$\Int$ સમક્રમી મંડળ છે તે સાબિત કરો.
\end{prob}

પરંતુ આપણું કામ પૂરું થયું નથી. $\Int$ પર સરવાળો અને ગુણાકાર
વ્યાખ્યાયિત કરવા ઉપરાંત આપણે ક્રમસંબંધ $\leq$ પણ વ્યાખ્યાયિત કર્યો
હતો અને તેનો વ્યવહાર પણ જેવો હોવો જોઈએ તેવો છે એ ચકાસવું પડે. વધુ
ચોક્કસ રીતે, $\Int$ \emph{ક્રમિત} મંડળ બનાવે છે એવું બતાવવું પડશે.\footnote{
	\olref[sfr][rel][ord]{def:linearorder}માંથી યાદ કરો કે સંપૂર્ણ ક્રમ
	સ્વવાચક, પરંપરિત, વિસંમિત અને તુલનાયુક્ત સંબંધ છે. ક્રમસંબંધોના
	સંદર્ભમાં તુલનાયુક્તતાને ક્યારેક \emph{ત્રિવિકલ્પતા} કહે છે, કારણ કે
	કોઈ પણ $a$ અને $b$ માટે $a \leq b \lor a = b \lor a \geq b$.}

\begin{defn}
\emph{ક્રમિત મંડળ} એવું સમક્રમી મંડળ છે જેમાં સંપૂર્ણ ક્રમસંબંધ
$\leq$ પણ હોય અને નીચેની શરતો સંતોષાય:
\begin{align*}
	a \leq b &\lif a + c \leq b + c\\
	(a \leq b \land 0 \leq c) &\lif a \times c \leq b \times c
\end{align*}
\end{defn}

\begin{prob}
$\Int$ ક્રમિત મંડળ છે તે સાબિત કરો.
\end{prob}

અગાઉની જેમ રચાયેલું~$\Int$ ક્રમિત મંડળ છે એવું બતાવવું શ્રમસાધ્ય
પણ નિયમિત કામ છે. તે તમારા માટે છોડી દઈએ છીએ.

આમ પૂર્ણાંકોનું કામ પૂરું થયું. હવે સંમેય સંખ્યાઓ માટે આવી જ બાબતો
બતાવવી છે. ખાસ કરીને $+$, $\times$ અને $\leq$ની આપેલી વ્યાખ્યાઓ તળે
સંમેય સંખ્યાઓ ક્રમિત \emph{ક્ષેત્ર} બનાવે છે એવું બતાવવું છે:
\begin{defn}\ollabel{orderedfield}
\emph{ક્રમિત ક્ષેત્ર} એવું ક્રમિત મંડળ છે જે આ વધારાની શરત પણ સંતોષે:
\begin{align*}
	\emph{ગુણાકારી વ્યસ્ત}& & (\forall a \in S \setminus \{0\})(\exists b \in S) a\times b& = 1
\end{align*}
\end{defn}

$\Int$ ક્રમિત મંડળ બનાવે છે એવું બતાવી લીધા પછી $\Rat$ ક્રમિત ક્ષેત્ર
બનાવે છે એવું બતાવવું સરળ, પણ શ્રમસાધ્ય છે.

\begin{prob}
$\Rat$ ક્રમિત ક્ષેત્ર છે તે સાબિત કરો.
\end{prob}

પૂર્ણાંકો અને સંમેય સંખ્યાઓનું કામ પૂરું કર્યા પછી માત્ર વાસ્તવિક
સંખ્યાઓ બાકી રહે છે. ખાસ કરીને, $\Real$ \emph{પૂર્ણ} ક્રમિત ક્ષેત્ર,
એટલે કે પૂર્ણતા ગુણધર્મ ધરાવતું ક્રમિત ક્ષેત્ર, બનાવે છે એવું બતાવવું
પડે. \olref[cuts]{realcompleteness}માં $\Real$ પૂર્ણતા ગુણધર્મ ધરાવે
છે એવું સ્થાપ્યું હતું. પરંતુ $\Real$ ક્રમિત ક્ષેત્ર છે તેની
(કંટાળાજનક) ચકાસણીની વિગતો હજી બાકી છે.

આ શ્રમસાધ્ય કસરત તરફ વળતાં પહેલાં થોડી વધુ ``તાત્કાલિક'' બાબતો
ચકાસવી પડશે. દાખલા તરીકે, કોઈ પણ કાપો $\alpha$ અને $\beta$ માટે
વ્યાખ્યાયિત કરેલો $\alpha + \beta$ ખરેખર \emph{કાપ} જ છે તેની ખાતરી
કરવી પડશે. તેની સાબિતી આ રહી:

\begin{proof}
$\alpha$ અને $\beta$ બંને કાપ હોવાથી $\alpha + \beta = \Setabs{p
+ q}{p \in \alpha \land q \in \beta}$ એ $\Rat$નો અરિક્ત ઉચિત ઉપગણ
છે. હવે કોઈ $p \in \alpha$ અને $q \in \beta$ માટે $x < p + q$ ધારો.
ત્યારે $x - p < q$, તેથી $x - p \in \beta$, અને $x = p + (x - p) \in
\alpha + \beta$. તેથી $\alpha + \beta$ એ $\Rat$નો આરંભિક ખંડ છે.
અંતે, કોઈ પણ $p + q \in \alpha + \beta$ માટે, $\alpha$ અને $\beta$
બંને કાપ હોવાથી એવા $p_1 \in \alpha$ અને $q_1 \in \beta$ છે કે
$p < p_1$ અને $q < q_1$; તેથી $p + q < p_1 + q_1 \in \alpha +
\beta$; આમ $\alpha + \beta$ને કોઈ મહત્તમ ઘટક નથી.
\end{proof}

આવી જ મહેનતથી તમે ચકાસી શકશો કે $\alpha - \beta$, $\alpha \times
\beta$ અને $\alpha \div \beta$ પણ કાપ છે (છેલ્લા કિસ્સામાં $\beta$
શૂન્ય-કાપ હોય તે કિસ્સો અવગણવો). એ પણ આપણે તમારા માટે જ છોડી દઈએ છીએ.

\begin{prob}
$\Real$ ક્રમિત ક્ષેત્ર છે તે સાબિત કરો.
\end{prob}

હવે એક નાની અધૂરી વિગત પૂરી કરીએ. \olref[cuts]{sec}માં આપણે સૂચવ્યું
હતું કે $\sqrt{2} = \Setabs{p \in \Rat}{p < 0 \text{ અથવા }p^2 < 2}$
લઈ શકાય. પરંતુ આ ગણ \emph{કાપ} છે એવું બતાવવું પડે. તેની સાબિતી આ રહી:

\begin{proof}
સ્પષ્ટપણે આ સંમેય સંખ્યાઓનો અરિક્ત ઉચિત આરંભિક ખંડ છે; તેથી તેને
મહત્તમ ઘટક નથી એટલું બતાવવું પૂરતું છે. ખાસ કરીને, $p^2 < 2$ ધરાવતી
ધન સંમેય સંખ્યા $p$ અને $q = \frac{2p+2}{p+2}$ હોય ત્યારે $p < q$ અને
$q^2 < 2$ બંને બતાવવાં પૂરતાં છે. $p < q$ જોવા માટે માત્ર નોંધો:
\begin{align*}
	p^2 &< 2\\
	p^2 + 2p &< 2 + 2p\\
	p(p + 2) &< 2 + 2p\\
	p &< \tfrac{2+2p}{p+2} = q
\end{align*}
$q^2 < 2$ જોવા માટે માત્ર નોંધો:
\begin{align*}
	p^2 &< 2\\
	2p^2 + 4p + 2 &< p^2 + 4p+ 4\\
	4p^2 + 8p + 4 &< 2(p^2 + 4p + 4)\\
	(2p+2)^2 & <2(p+2)^2\\
	\tfrac{(2p+2)^2}{(p+2)^2} &< 2\\
	q^2 &< 2
\end{align*}
\end{proof}
\end{document}
